% Generated draft: requires mathematical and source review.
\documentclass[11pt]{article}
\usepackage[margin=25mm]{geometry}
\usepackage{fontspec}
\setmainfont{Latin Modern Roman}
\usepackage{amsmath,amssymb,mathtools,unicode-math}
\setmathfont{Latin Modern Math}
\usepackage{enumitem,needspace,xurl,hyperref}
\setlength{\parindent}{0pt}
\setlength{\parskip}{5pt}
\setlength{\emergencystretch}{3em}
\newcommand{\N}{\mathbb{N}}
\newcommand{\Z}{\mathbb{Z}}
\newcommand{\Q}{\mathbb{Q}}
\newcommand{\R}{\mathbb{R}}
\newcommand{\C}{\mathbb{C}}
\newcommand{\F}{\mathbb{F}}
\newcommand{\A}{\mathbb{A}}
\newcommand{\PP}{\mathbb{P}}
\newcommand{\E}{\mathbb{E}}
\newcommand{\eps}{\varepsilon}
\newcommand{\dd}{\,\mathrm{d}}
\begin{document}
\section*{UnsolvedMath ID: 30000519}
\textbf{Source problem code:} catalogue entry
\subsection{Definitions and mathematical statement}
\textbf{Three Stokes-phenomenon tasks.} For a \((1+1)\)-dimensional PDE, determine the possible exponential phases, the regions where each Stokes line is active, and the multiplier switched on when an active Stokes line is crossed. The variables conventionally represent one spatial coordinate \(x\) and one time coordinate \(t\). The PDE, small- or large-parameter regime, distinguished solution, and initial or boundary data are not supplied, so these are model-dependent research tasks rather than one specified theorem.\par In a conditional small-parameter notation, exponential contributions may have the form \(A_j(x,t;\varepsilon)e^{-\chi_j(x,t)/\varepsilon}\), with \(\chi_j\) the phases or singulants and \(A_j\) asymptotic amplitudes. The first task is to determine which such phases actually arise for the given equation and solution, including whatever branch and interaction conditions it imposes. This schematic expression does not assume a particular PDE, a finite phase list, or a derived eikonal equation.\par The second task is to identify the relevant Stokes curves and distinguish geometrically possible curves from active portions along which a subdominant exponential contribution is switched on. Under a common convention, a candidate switching condition involves alignment of relative phases, such as \(\operatorname{Im}(\chi_j-\chi_k)=0\), together with the prescribed real-part dominance and orientation. The source does not fix Stokes versus anti-Stokes terminology, the parameter argument, or the real or complexified domain; no sign condition is chosen silently. The third task is to determine the corresponding coefficient or jump of the normalized subdominant contribution when an active curve is crossed. That multiplier depends on normalization, selected solution, continuation path, and data. It is not automatically one universal constant. A Stokes jump in an asymptotic representation is not itself a discontinuity of the exact solution.\par The supplied triage checked 2026-08-21 labels the item open but explicitly calls the three questions a general research program, with phases, activity, and multipliers depending on the particular differential equation. It reports no universal resolution from the extracted statement. This dated qualification is retained rather than promoted to an unformulated universal conjecture.
\subsection{Short English statement}
For a specified PDE asymptotic problem, identify its exponential phases, determine where switching actually occurs, and compute the normalized multiplier. The equation and asymptotic conventions are missing, so no universal phase list, active-curve rule, or multiplier is asserted.
\subsection{Sources}
\begin{itemize}\item Local UnsolvedMath registry, permanent ID 30000519, code catalogue entry; \emph{Stokes Phases and Multipliers for Nonlinear PDEs}, Oberwolfach Reports Open Problems.\item \emph{Applications of Asymptotic Analysis (2007)}, Oberwolfach Reports, DOI 10.4171/owr/2006/28, pp. 1663--1730. \url{https://doi.org/10.4171/owr/2006/28}. Proposed year and year tag: 2006; no proposer named.\item Supplied literature label \emph{original Oberwolfach report}: \url{https://doi.org/10.4171/owr/2006/28}.\item Supplied literature label \emph{Stokes phenomenon survey}: \url{https://doi.org/10.1137/1031090}.\item Literature triage checked 2026-08-21, status \texttt{open} with a research-program qualification; record updated 2026-09-16. The \texttt{corrected\_verified} field reports source checking, not independently verified here.\end{itemize}
\end{document}
